\chapter{Probabilistic Foundations: Arrivals, Service, and Little's Law}\label{sec:L1}

This lecture defines the vocabulary of queueing theory and proves three tools that
every later lecture relies on: the memorylessness of the exponential distribution,
the Poisson process together with PASTA, and Little's law. We assume undergraduate
probability (conditional expectation, the strong law of large numbers, probability
generating functions).

\section{Queueing systems and performance measures}

\begin{definition}[Queueing system]\label{L1:def:system}
A \term{queueing system} consists of four elements.
\begin{enumerate}
\item \term{Customers} (or jobs): entities that arrive at the system, require some
  amount of processing, and leave. The arrival time of the $k$-th customer is $a_k$,
  numbered so that $0\le a_1\le a_2\le\cdots$.
\item The \term{arrival process}: the sequence of arrival times $(a_k)_{k\ge1}$, or
  equivalently the counting process $A(t)=\#\{k: a_k\le t\}$.
\item A \term{server} and \term{service times}: the server is the resource that
  processes customers. Customer $k$ needs time $S_k$ when it has the server to
  itself. $S_k$ is also called the \term{work} or service requirement of customer $k$.
\item A \term{waiting room} and a \term{service discipline}: the space where customers
  arriving at a busy server wait, and the rule that decides whom to serve, when, and
  how much.
\end{enumerate}
\end{definition}

\begin{definition}[Service disciplines]\label{L1:def:discipline}
\term{FIFO} (first in, first out) serves one customer at a time, to completion, in
order of arrival. \term{Processor sharing} (PS) serves all customers in the system
simultaneously at equal rates; it is defined formally in \cref{sec:L2}.
A discipline under which the server never idles while work is present is called
\term{work-conserving}. Both FIFO and PS are work-conserving.
\end{definition}

\begin{definition}[Kendall notation]\label{L1:def:kendall}
A queueing system is written $A/S/c/K/N$. Here $A$ is the distribution of the
interarrival times, $S$ the distribution of the service times, $c$ the number of
servers, $K$ the system capacity ($\infty$ if omitted), and $N$ the size of the
customer population ($\infty$ if omitted). The distribution symbols are $M$ (Markov:
exponential and independent), $D$ (deterministic: constant), and $G$ (general:
an arbitrary distribution, i.i.d.).
\end{definition}

For example, M/M/1 has Poisson arrivals, exponential service and one server
(\cref{sec:L2}); M/G/1 allows a general service distribution (\cref{sec:L3});
M/M/1//N is a closed system with only $N$ customers (\cref{sec:L4}); and
M/G/$\infty$ is an infinite-server system with no waiting (\cref{sec:L4}).

\begin{definition}[Performance measures]\label{L1:def:measures}
Let $d_k$ be the time at which customer $k$ leaves the system.
\begin{enumerate}
\item The \term{sojourn time} (response time) is $W_k=d_k-a_k$.
\item The \term{waiting time} $W_{q,k}$ is the time from the arrival of customer $k$
  until its service first begins. For a FIFO single server, $W_k=W_{q,k}+S_k$.
\item The \term{number in system} is $N(t)=\#\{k: a_k\le t<d_k\}$. The number in the
  waiting room is written $N_q(t)$.
\end{enumerate}
\end{definition}

Quantities in queueing theory are defined by two different kinds of averages, and the
distinction will be essential several times.

\begin{definition}[Time averages and customer averages]\label{L1:def:averages}
The \term{time average} of the process $N(t)$ is
$L=\lim_{t\to\infty}\frac1t\int_0^t N(s)\dd s$, and the \term{customer average} of the
sequence $(W_k)$ is $W=\lim_{n\to\infty}\frac1n\sum_{k=1}^n W_k$ (when the limits
exist). $L_q$ and $W_q$ are defined in the same way. If the limits are constants with
probability one, then in the stationary regime they equal $\E[N]$ and $\E[W]$
respectively.
\end{definition}

\begin{definition}[Arrival rate, load, stability]\label{def:load}
The \term{arrival rate} is $\lambda=\lim_{t\to\infty}A(t)/t$.
The \term{load} (utilisation) of a single-server system with mean service time
$\E[S]$ is
\[
  \rho=\lambda\,\E[S].
\]
$\rho$ is the amount of work arriving per unit time and is dimensionless.
The system is \term{stable} if $N(t)$ has a proper limiting distribution (or at least
$L<\infty$) and every customer leaves in finite time.
\end{definition}

A single server processes at most one unit of work per unit time, so if $\rho>1$ the
unfinished work grows linearly and the system cannot be stable. For a work-conserving
G/G/1 system (i.i.d.\ interarrival and service times) it is known that $\rho<1$ is
essentially necessary and sufficient for stability \cite{loynes1962}. For the models
in these notes we assume $\rho<1$ as the stability condition and verify it model by
model.

\section{From a serving system to a queueing model}\label{L1:sec:model}

A queueing model of a serving system keeps what decides which request waits and for
how long, and leaves out the rest. \Cref{fig:deployment} shows the model of a
disaggregated deployment used throughout these notes. The queueing objects it uses
(multi-server queues, processor sharing, delay stations, closed loops) are defined and
analysed in later lectures; here the serving system itself is defined, first in the
picture and then exactly, so that every later theorem has a place in it.

\begin{figure}[t]
\centering
\begin{tikzpicture}[x=0.8cm, y=0.8cm, >=Stealth, font=\small, line width=0.6pt,
  station/.style={circle, draw, thick, minimum size=10mm, align=center, inner sep=1pt},
  note/.style={font=\footnotesize, align=center, text=black!75}]
  % ---------- prefill instance ----------
  \foreach \i in {0,1}
    \draw[fill=orange!10] (0.4+0.5*\i,-0.3) rectangle ++(0.45,0.6);
  \draw (0.35,-0.3) -- (1.4,-0.3) (0.35,0.3) -- (1.4,0.3);
  \node[draw, thick, rounded corners=1pt, fill=orange!10, minimum width=12mm, minimum height=9mm] (SP) at (2.5,0) {};
  \foreach \x in {-0.36,0,0.36}
    \foreach \y in {-0.17,0.17}
      \draw[orange!60!black, fill=white] ($(SP)+(\x-0.13,\y-0.11)$) rectangle ++(0.26,0.22);
  \node[station, fill=blue!8] (P) at (4.4,0) {FIFO};
  \draw[->] (SP.east) -- (P.west);
  \draw[fill=gray!15, draw=gray!60] (0.2,-1.4) rectangle (5.3,-0.85);
  \node[font=\scriptsize, text=black!75] at (2.75,-1.12) {prefix cache of paused sessions};
  \node[note] at (2.5,0.85) {$m_P$ slots};
  \node[note] at (4.4,0.9) {prefill};
  \node[draw, dashed, rounded corners=3mm, inner sep=2.5mm, fit={(0.2,-1.4) (5.3,1.2)}] (MP) {};
  \node[note, anchor=south] at (MP.north) {\textbf{Prefill instance} (pool $M_P$)};
  % ---------- link ----------
  \node[station, fill=green!8] (L) at (7.0,0) {PS};
  \node[note, anchor=north] at (7.0,-0.65) {KV link\\bandwidth $B$};
  % ---------- decode instance ----------
  \foreach \i in {0,1}
    \draw[fill=orange!10] (8.9+0.5*\i,-0.3) rectangle ++(0.45,0.6);
  \draw (8.85,-0.3) -- (9.9,-0.3) (8.85,0.3) -- (9.9,0.3);
  \node[draw, thick, rounded corners=1pt, fill=orange!10, minimum width=12mm, minimum height=9mm] (SD) at (11.0,0) {};
  \foreach \x in {-0.36,0,0.36}
    \foreach \y in {-0.17,0.17}
      \draw[orange!60!black, fill=white] ($(SD)+(\x-0.13,\y-0.11)$) rectangle ++(0.26,0.22);
  \node[station, fill=red!8] (D) at (12.9,0) {PS};
  \draw[->] (SD.east) -- (D.west);
  \node[note] at (11.0,0.85) {$m_D$ slots};
  \node[note] at (12.9,-0.95) {batch $\varphi(m)$};
  \node[draw, dashed, rounded corners=3mm, inner sep=2.5mm, fit={(8.8,-1.25) (14.0,1.2)}] (MD) {};
  \node[note, anchor=south] at (MD.north) {\textbf{Decode instance} (pool $M_D$)};
  % ---------- tool call ----------
  \node[draw, thick, rounded corners, fill=teal!8, minimum width=20mm, minimum height=9mm] (T) at (7.0,-3.1) {};
  \foreach \k in {-2,-1,0,1,2}
    \draw[teal!60!black] ($(T)+(0.36*\k,0)$) circle (0.13);
  \node[note, anchor=north] at ($(T.south)+(0,-0.05)$) {\textbf{Tool call}: delay station};
  % ---------- flows ----------
  \draw[->] (-3.0,0) -- (0.3,0);
  \node[note, anchor=south] at (-2.1,0.05) {new sessions\\Poisson $\Lambda$};
  \draw[->] (P.east) -- (L.west);
  \draw[->] (L.east) -- (8.8,0);
  \draw[->] (D.east) -- (15.4,0) node[note, right] {ends};
  \draw[->] (D.south) |- node[note, right, pos=0.25] {resume w.p.\ $p$} (T.east);
  \draw[->] (T.west) -| (-0.2,0) -- (0.3,0);
  \node[note, anchor=east] at (-0.3,-2.2) {next turn\\(hit or miss)};
\end{tikzpicture}
\caption{A disaggregated deployment as a queueing model. A turn waits for KV memory
at a prefill instance (it is admitted only when its whole context fits, in arrival
order), is prefilled one request at a time in arrival order, sends its KV cache over a
link whose bandwidth concurrent transfers share, waits for KV memory at a decode
instance, and decodes in a batch that shares the instance's throughput. Sessions loop
through a tool call back to the prefill instance, where the prefix cache decides
whether the next turn is a hit (short prefill work) or a miss (the whole context is
recomputed).}
\label{fig:deployment}
\end{figure}

The figure is a picture, not a definition: it does not say when a waiting request is
admitted, what happens to its KV memory afterwards, or who decides a hit. The rest of
this section makes it a definition. We first define a small language in which a
serving deployment is a program, and give it a precise meaning
(\cref{L1:sec:language}); the language knows nothing about tokens or prefill costs.
We then write the deployment of \cref{fig:deployment} as one program, with its
workload and costs as the program's parameters (\cref{L1:sec:replica}), and read off
what that program implies (\cref{L1:sec:implies}). The tools of later lectures are
needed only to compute averages of the quantities defined here.

\subsection{A language for serving deployments}\label{L1:sec:language}

\paragraph{Requests and sessions.} A \emph{request} is one turn of an agent: it
arrives at a prefill instance, waits, is prefilled, transferred and decoded, and
leaves. A \emph{session} is the agent program that sends these requests: it returns
after each tool call and leaves only when the program ends. New sessions arrive from
outside as a Poisson process, but the requests of a live session do not arrive
independently of the system: the next one exists only after the previous one has been
decoded and its tool call has returned. In the language the object that moves is
therefore the session, and each request is one pass of the session through a loop.

A serving deployment does four things to a request: it makes it wait until its KV
cache fits in a memory pool, runs it on a stage (a prefill instance, a link, a decode
instance, a tool), frees its memory (possibly keeping the prefix cached), and sends it
somewhere next. We give each of these an instruction. The resulting language, which we
call \textsc{Route}, describes a system by a \emph{deployment} (its memory pools and
stages), a \emph{workload} (how sessions arrive and how their turns evolve), and a
\emph{route}, the path a session follows through the deployment.

\begin{definition}[Syntax of \textsc{Route}]\label{L1:def:syntax}
A \textsc{Route} program is a triple $(\mathcal D,\mathcal W,\pi_0)$ of a deployment, a
workload and a route.
\begin{align*}
  \text{deployment item}\quad d \;&::=\; \kw{pool}\ m\ (M,\,E)
    \;\mid\; \kw{stage}\ s : k\\
  \text{stage kind}\quad k \;&::=\; \mathrm{serial} \;\mid\; \mathrm{shared}(\varphi)
    \;\mid\; \mathrm{external} \;\mid\; \mathrm{step}(B,\tau)\\
  \text{workload}\quad \mathcal W \;&::=\; \kw{arrivals}\ (\Lambda,\,y_0,\,G)\\
  \text{action}\quad \alpha \;&::=\; \kw{turn}
    \;\mid\; \kw{admit}\ m\ c
    \;\mid\; \kw{free}\ m
    \;\mid\; \kw{free}\ m\ \kw{cache}\ \ell\\
    &\phantom{::=}\;\;\mid\; \kw{run}\ s\ w
    \;\mid\; \kw{run}\ s\ (w_{\mathrm p}, w_{\mathrm d})\ \kw{growing}\ m\\
  \text{route}\quad \pi \;&::=\; \epsilon
    \;\mid\; \alpha;\pi
    \;\mid\; \kw{loop}\ \pi
    \;\mid\; \kw{branch}_p\ \pi_1\ \pi_2
    \;\mid\; \kw{end}
\end{align*}
Here $m$ names a KV memory pool of $M$ bytes with eviction order $E$, $s$ names a
stage, and $p\in[0,1]$. A \emph{serial} stage runs one request at a time, in arrival
order, at rate $1$ (a prefill instance); a \emph{shared}$(\varphi)$ stage runs all its
requests at once and splits a throughput $\varphi(n)$ equally among the $n$ present
(a decode batch, a link); an \emph{external} stage runs every request on its own, with
no waiting (a tool call); a \emph{step}$(B,\tau)$ stage is an engine that
runs iterations, handing a budget of $B$ tokens per iteration to its jobs
in their order of entry, one token to a job in its decode phase and as many
as its remaining prefill work allows to a job in its prefill phase, and
taking $\tau(\cdot)$ seconds for an iteration of the tokens it scheduled
(a colocated replica with chunked prefill). The workload has a session arrival rate $\Lambda$, a
measurable \term{attribute space} $\mathcal X$ with the attributes $y_0\in\mathcal X$ of
a new session, and a \term{turn kernel} $G$, a transition kernel on $\mathcal X$:
$G(y,\cdot)$ is the law of a session's next-turn attributes given its current ones. A
session's attributes are $x=(y,H)\in\mathcal X\times\{0,1\}$, where $y$ is set by the
workload and the \term{hit indicator} $H$ by the deployment. The arguments $c$ (bytes to
allocate), $\ell$ (bytes to keep cached, $\ell\le c$) and $w$ (work) are nonnegative
measurable functions of $x$. Sequencing is associative, with $\epsilon$ as unit.
\end{definition}

We write $\kw{hold}\ m\ c\ \{\pi\}\ \kw{cache}\ \ell$ for
$\kw{admit}\ m\ c;\ \pi;\ \kw{free}\ m\ \kw{cache}\ \ell$: every program we write
frees what it admits, and the scoped form makes that balance syntactic.
Informally: $\kw{turn}$ draws the session's next-turn attributes from $G$;
$\kw{admit}\ m\ c$ joins pool $m$'s admission queue and leaves it with $c$ bytes
allocated; $\kw{free}\ m\ \kw{cache}\ \ell$ frees the allocation but keeps $\ell$ bytes
of it as a cached prefix, which the pool may evict later; $\kw{run}\ s\ w$ joins stage
$s$ with $w$ units of work and continues when the work is done;
$\kw{run}\ s\ (w_{\mathrm p},w_{\mathrm d})\ \kw{growing}\ m$ joins a step stage with
$w_{\mathrm p}$ prefill and $w_{\mathrm d}$ decode tokens, and before each iteration
enlarges its allocation in pool $m$ to the tokens it will have computed; $\kw{branch}_p$
continues with $\pi_1$ with probability $p$ and with $\pi_2$ otherwise; $\kw{end}$ ends
the session.

\paragraph{Semantics.} The meaning of a program is a stochastic process whose state is
the following configuration. For a finite sequence $q$ we write $r\cdot q$ for $q$
with $r$ in front, $q\cdot r$ for $q$ with $r$ appended, $q-r$ for $q$ with $r$ removed,
and for a partial map $f$ into $\R_+$ (or into pairs whose first entry is in $\R_+$) we
write $|f|$ for the sum of its (first) values.

\begin{definition}[Configuration]\label{L1:def:config}
A \term{configuration} $\Gamma$ of a program $(\mathcal D,\mathcal W,\pi_0)$ consists of
\begin{enumerate}[label=(\alph*)]
\item the time $t\ge0$ and the finite set $\mathcal R$ of live sessions;
\item for each session $r\in\mathcal R$: its attributes $x_r=(y_r,H_r)$, its
  continuation $\pi_r$ (the rest of its route), its status
  $\mathrm{st}_r\in\{\mathsf{ready}\}\cup\{\mathsf{queued}(m)\}\cup\{\mathsf{in}(s)\}$,
  and its remaining work $\omega_r\ge0$;
\item for each pool $m$: its admission queue $q_m$, a finite sequence of sessions; its
  allocation $h_m$, a partial map from sessions to bytes; and its prefix cache $C_m$, a
  partial map from sessions to pairs (bytes, time of last use);
\item for each stage $s$: the sequence $Q_s$ of sessions at the stage, in order of
  entry.
\end{enumerate}
Every reachable configuration satisfies the memory invariant
\begin{equation}\label{L1:eq:invariant}
  |h_m|+|C_m|\le M\qquad\text{for every pool } m \text{ of } M \text{ bytes}.
\end{equation}
\end{definition}

The program runs by two kinds of steps. \emph{Commands} take no time and change the
configuration when their guard holds. \emph{Flow} lets time pass while no command is
enabled. For a pool with eviction order $E$, $\mathrm{evict}_E(C,b)$ deletes entries of
$C$ in the order $E$ until $|C|\le b$; for LRU the order is increasing time of last use.
We write ``$r$ at $\alpha;\pi'$'' for $\mathrm{st}_r=\mathsf{ready}$ and
$\pi_r=\alpha;\pi'$.

\begin{definition}[Operational semantics of \textsc{Route}]\label{L1:def:semantics}
The commands are the following guarded assignments, written
$[\textsc{name}]\ \text{guard}\ \longrightarrow\ \text{effect}$.
\[
\renewcommand{\arraystretch}{1.25}
\begin{array}{@{}l@{\ }l@{\ }l@{}}
{[\textsc{Turn}]} & r \text{ at } \kw{turn};\pi'
  &\longrightarrow\ \text{draw } y'\sim G(y_r,\cdot);\ y_r:=y';\ \pi_r:=\pi'\\
{[\textsc{Queue}]} & r \text{ at } \kw{admit}\ m\ c;\pi'
  &\longrightarrow\ q_m:=q_m\cdot r;\ \mathrm{st}_r:=\mathsf{queued}(m);\ \pi_r:=\pi'\\
{[\textsc{Admit}]} & q_m=r\cdot q',\ |h_m|+c(x_r)\le M
  &\longrightarrow\ H_r:=\1\{r\in\operatorname{dom}C_m\};\ C_m:=C_m-r;\\
&&\phantom{\longrightarrow}\ C_m:=\mathrm{evict}_E\bigl(C_m,\,M-|h_m|-c(x_r)\bigr);\\
&&\phantom{\longrightarrow}\ h_m(r):=c(x_r);\ q_m:=q';\ \mathrm{st}_r:=\mathsf{ready}\\
{[\textsc{Free}]} & r \text{ at } \kw{free}\ m\ [\kw{cache}\ \ell];\pi'
  &\longrightarrow\ h_m:=h_m-r;\ [\,C_m(r):=(\ell(x_r),t)\,];\\
&&\phantom{\longrightarrow}\ \pi_r:=\pi'\\
{[\textsc{Run}]} & r \text{ at } \kw{run}\ s\ w;\pi'
  &\longrightarrow\ \omega_r:=w(x_r);\ Q_s:=Q_s\cdot r;\\
&&\phantom{\longrightarrow}\ \mathrm{st}_r:=\mathsf{in}(s);\ \pi_r:=\pi'\\
{[\textsc{Finish}]} & \mathrm{st}_r=\mathsf{in}(s),\ \omega_r=0
  &\longrightarrow\ Q_s:=Q_s-r;\ \mathrm{st}_r:=\mathsf{ready}\\
{[\textsc{Branch}]} & r \text{ at } \kw{branch}_p\ \pi_1\ \pi_2;\pi'
  &\longrightarrow\ \pi_r:=\pi_1;\pi' \text{ w.p.\ } p,\ \text{else } \pi_r:=\pi_2;\pi'\\
{[\textsc{Loop}]} & r \text{ at } \kw{loop}\ \pi;\pi'
  &\longrightarrow\ \pi_r:=\pi;\kw{loop}\ \pi;\pi'\\
{[\textsc{End}]} & r \text{ at } \kw{end};\pi'
  &\longrightarrow\ \mathcal R:=\mathcal R-r
\end{array}
\]
The brackets in \textsc{Free} are present exactly when the action has $\kw{cache}$.
All random draws are independent of each other and of the past. Commands are
\emph{urgent}: while any command is enabled, one is executed (the enabled command of
the session that arrived first). When none is enabled, \textsc{Flow} lets time
advance: each session at a stage works off its remaining work at rate
\begin{equation}\label{L1:eq:rates}
  v_r=\begin{cases}
    \1\{r \text{ is the first entry of } Q_s\} & s:\mathrm{serial},\\[0.2em]
    \varphi(|Q_s|)/|Q_s| & s:\mathrm{shared}(\varphi),\\[0.2em]
    1 & s:\mathrm{external},
  \end{cases}
  \qquad \frac{\dd\omega_r}{\dd t}=-v_r,
\end{equation}
until the first time some $\omega_r$ reaches $0$ (which enables \textsc{Finish}) or a
session arrives. At each epoch of the Poisson process of rate $\Lambda$, the command
\textsc{Arrive} adds a new session $r$ to $\mathcal R$ with $y_r=y_0$, $H_r=0$,
$\pi_r=\pi_0$ and $\mathrm{st}_r=\mathsf{ready}$.
\end{definition}

A step stage does not flow continuously: it advances by iterations. At
the start of an iteration its jobs are taken in order of entry and each is
handed $\min(\text{remaining},\ \text{budget left})$ tokens (one for a decode
phase); a $\kw{growing}$ job first enlarges its allocation, and if the pool
has no room the most recently admitted holder is \emph{preempted}: its job
leaves the stage, its allocation is freed with the prefix it computed
cached, and it rejoins the head of the pool's queue to be admitted again.
The iteration ends $\tau$ seconds later and its tokens are then applied.

\textsc{Admit} is where the admission rule lives. Only the first session in the queue
can be admitted, so a request that does not fit blocks every request behind it; the
guard counts only the bytes allocated to admitted requests, not the cache, so a cached
prefix never blocks an admission; instead, the eviction in the second line of
\textsc{Admit} deletes as many cached prefixes as the new allocation needs, which keeps
\eqref{L1:eq:invariant}. The hit indicator is decided at the same moment, by whether
the session's own prefix survived until its admission. \textsc{Free} also keeps the
invariant, since it frees $c$ bytes and caches $\ell\le c$ of them.

\paragraph{Production schedulers.} Three refinements of \textsc{Admit} are
needed to write a production engine such as vLLM as a program, and none
changes the invariant. (i) The queue of a pool may be served by a step
stage: admission happens only at the start of an iteration, with the budget
the running jobs leave, so a waiting request's cached prefix is not
protected until then and can be evicted while it waits. (ii) The guard may
ask for more than is allocated: vLLM admits a request only if its whole
prompt fits, but allocates only the blocks of its first chunk and grows
from there. (iii) A request may reuse only part of its cached prefix (the
blocks it shares with its previous turn); the rest stays cached, unusable,
until it is evicted. With these rules the program of the vLLM v1 engine
reproduces the real scheduler step for step, and it keeps the cached
prefixes of finished sessions, as \textsc{End} does: the cache does not
know that a session has left.

The semantics is well defined for programs in which every $\kw{loop}$ body contains a
$\kw{run}$ with positive work (a step stage's jobs are not in \eqref{L1:eq:rates}: they
advance only at the end of an iteration): then between two flow steps each session executes only
finitely many commands. It is also an algorithm. \Cref{alg:route} runs it, and a
discrete-event simulator of a serving system is an implementation of this loop.

\begin{algorithm}[t]
\caption{Interpreter of a \textsc{Route} program up to time $t_{\max}$}\label{alg:route}
\begin{algorithmic}[1]
\Require deployment $\mathcal D$, workload $\mathcal W$, route $\pi_0$, horizon $t_{\max}$
\State $t\gets0$;\ $\mathcal R\gets\emptyset$;\ all queues, allocations, caches and stages empty
\State $a\gets$ a draw from $\mathrm{Exp}(\Lambda)$ \Comment{next session arrival}
\While{$t<t_{\max}$}
  \While{some command is enabled} \Comment{zero-time phase}
    \State execute the enabled command of the earliest-arrived session
  \EndWhile
  \State compute $v_r$ for every session at a stage by \eqref{L1:eq:rates}
  \State $e\gets$ the earliest end of a running iteration of a step stage ($\infty$ if none)
  \State $\Delta\gets\min\bigl(a-t,\ e-t,\ \min_{r:\,v_r>0}\omega_r/v_r\bigr)$
  \For{every session $r$ at a stage}
    \State $\omega_r\gets\omega_r-v_r\Delta$
  \EndFor
  \State $t\gets t+\Delta$
  \If{$t=a$}
    \State \textsc{Arrive};\ $a\gets t+{}$a draw from $\mathrm{Exp}(\Lambda)$
  \EndIf
  \If{$t=e$} \Comment{apply the iteration's tokens; the next one is assigned}
    \State end the iteration; start the next one after the zero-time phase
  \EndIf
\EndWhile
\end{algorithmic}
\end{algorithm}

\subsection{The disaggregated replica}\label{L1:sec:replica}

A program needs a workload and the functions $c$, $\ell$, $w$. For the deployment of
\cref{fig:deployment} they come from the agentic workload and from the costs of one
replica.

\begin{definition}[Agentic workload]\label{L1:def:workload}
Sessions arrive at the epochs $A_1<A_2<\cdots$ of a Poisson process of rate $\Lambda$
(\cref{def:poisson}). The attributes of a turn are $y=(K,n,o,T,Z)\in\mathcal X=\R_+^5$.
\begin{enumerate}[label=(\alph*)]
\item \emph{Tokens.} Turn $j$ of session $i$ carries a prefix of $K_{ij}$ tokens that
  the prefill instance has already processed, $n_{ij}$ new tokens, and produces
  $o_{ij}$ output tokens. Its context is $T_{ij}=K_{ij}+n_{ij}$, and
  \begin{equation}\label{L1:eq:context}
    K_{i1}=0,\qquad K_{i,j+1}=T_{ij},\qquad n_{i,j+1}\ge o_{ij}.
  \end{equation}
  The new tokens of the next turn contain the previous output, which was produced on
  the decode instance, and the tool's output.
\item \emph{Think times.} After turn $j$ the session calls a tool for a time
  $Z_{ij}\ge0$ with mean $Z$.
\item \emph{Turn kernel.} A new session has $y_0=0$. Given the current turn $y$,
  $G(y,\cdot)$ draws $(n',o',Z')$ with $n'\ge o$ and sets $K'=T$ and $T'=K'+n'$, so that
  successive turns satisfy \eqref{L1:eq:context}.
\end{enumerate}
\end{definition}

\begin{definition}[Costs]\label{L1:def:costs}
A replica has prefill constants $a,b>0$, a KV size of $\kappa$ bytes per token, a link
of bandwidth $B$ with fixed overhead $x_0$, a decode capacity function $\varphi$, and
memory pools of $M_P$ and $M_D$ bytes. For turn $(i,j)$:
\begin{enumerate}[label=(\alph*)]
\item Prefilling $n$ new tokens on a cached prefix of $K$ tokens takes
  \begin{equation}\label{L1:eq:prefill}
    P(n,K)=a\,n+b\,n\,(K+n/2)
  \end{equation}
  seconds. With the hit indicator $H_{ij}$ (whether the prefix is still cached when
  the turn is admitted), the prefill work is
  \begin{equation}\label{L1:eq:work}
    S_{ij}=H_{ij}\,P(n_{ij},K_{ij})+(1-H_{ij})\,P(T_{ij},0)
    =P(n_{ij},K_{ij})+(1-H_{ij})\,\Delta S(K_{ij}),
  \end{equation}
  where $\Delta S(K)=aK+bK^2/2$ is the extra work of a miss.
\item The transfer of the context to the decode instance is the work
  $X_{ij}=x_0+\kappa\,T_{ij}/B$ seconds.
\item The decode work is $D_{ij}=o_{ij}$ tokens. A batch of $m$ running turns takes
  $\tau(m)$ seconds per step and advances each turn by one token, so its capacity is
  $\varphi(m)=m/\tau(m)$ tokens per second.
\item The memory footprint is $c_{ij}=\kappa\,T_{ij}$ bytes on either instance.
\end{enumerate}
\end{definition}

\begin{example}[The disaggregated replica]\label{L1:ex:program}
The deployment of \cref{fig:deployment} is the program
\begin{center}
\begin{tabular}{@{}l@{\qquad}l@{}}
$\kw{pool}\ \mathsf{mem}_P\ (M_P,\ \mathrm{LRU})$ &
$\kw{stage}\ \mathsf{prefill} : \mathrm{serial}$ \\
$\kw{pool}\ \mathsf{mem}_D\ (M_D,\ \mathrm{LRU})$ &
$\kw{stage}\ \mathsf{link} : \mathrm{shared}(1)$ \\
$\kw{arrivals}\ (\Lambda,\ y_0,\ G)$ &
$\kw{stage}\ \mathsf{decode} : \mathrm{shared}(\varphi)$ \\
&
$\kw{stage}\ \mathsf{tool} : \mathrm{external}$
\end{tabular}
\\[0.6em]
\begin{tabular}{@{}l@{\qquad}l@{}}
$\pi_0 \;=\; \kw{turn};\ \kw{loop}\ \bigl($ & \\
\quad $\kw{admit}\ \mathsf{mem}_P\ c;$ & \emph{wait for KV memory at the prefill instance}\\
\quad $\kw{run}\ \mathsf{prefill}\ S;$ & \emph{the first token is out when this ends}\\
\quad $\kw{run}\ \mathsf{link}\ X;$ & \emph{KV transfer, bandwidth shared}\\
\quad $\kw{free}\ \mathsf{mem}_P\ \kw{cache}\ \kappa T;$ & \emph{memory freed, prefix stays cached}\\
\quad $\kw{admit}\ \mathsf{mem}_D\ c;$ & \emph{wait for KV memory at the decode instance}\\
\quad $\kw{run}\ \mathsf{decode}\ D;$ & \emph{batched decode}\\
\quad $\kw{free}\ \mathsf{mem}_D;$ & \\
\quad $\kw{branch}_p\ (\kw{run}\ \mathsf{tool}\ Z;\ \kw{turn})\ (\kw{end})\ \bigr)$ &
  \emph{tool call and next turn, or end}
\end{tabular}
\end{center}
with the workload $(\Lambda,y_0,G)$ of \cref{L1:def:workload} and $c$, $S$, $X$, $D$
from \cref{L1:def:costs}. Because $\kw{branch}_p$ ends the loop with probability $1-p$
after each turn, session $i$ has $J_i$ turns with
$\Prob(J_i\ge j+1\mid J_i\ge j)=p$, so $\E[J_i]=1/(1-p)$. Other deployments are other
programs: a replica without a prefix cache drops $\kw{cache}\ \kappa T$, and every turn
is then a miss (\cref{L1:exr:nocache}); a colocated replica is a shorter program with
one pool and no link (\cref{L1:exr:colocated}).
\end{example}

\subsection{What the program implies}\label{L1:sec:implies}

Running the program of \cref{L1:ex:program} produces, for every request, seven epochs.
Number requests $k=1,2,\dots$ in the order of their \textsc{Queue} at $\mathsf{mem}_P$.
We write $\sigma(k)$ for the session of request $k$, and $S_k$, $X_k$, $D_k$, $c_k$,
$H_k$ and $Z_k$ for its quantities. Let
\[
  a_k\le\alpha_k\le s_k\le f_k\le e_k\le\alpha'_k\le d_k
\]
be the times of its \textsc{Queue} at $\mathsf{mem}_P$, its \textsc{Admit} there, the
moment it becomes the first entry of $Q_{\mathsf{prefill}}$, its \textsc{Finish} at
$\mathsf{prefill}$, its \textsc{Finish} at $\mathsf{link}$, its \textsc{Admit} at
$\mathsf{mem}_D$, and its \textsc{Finish} at $\mathsf{decode}$. Let
$R_P(t)=\{k:\alpha_k\le t<e_k\}$ be the requests holding memory at the prefill
instance.

\begin{lemma}[Sample-path equations of the replica]\label{L1:lem:paths}
Under \cref{L1:def:semantics}, the program of \cref{L1:ex:program} satisfies
\begin{align}
  \alpha_k&=\inf\Bigl\{t\ge\max(a_k,\alpha_{k-1}):
      \textstyle\sum_{l\in R_P(t)}c_l+c_k\le M_P\Bigr\},\label{L1:eq:admitP}\\
  s_k&=\max(\alpha_k,f_{k-1}),\qquad f_k=s_k+S_k,\label{L1:eq:lindley}\\
  a_{k^+}&=d_k+Z_k\quad\text{if the session continues},\label{L1:eq:feedback}
\end{align}
where $k^+$ is the next request of session $\sigma(k)$, and
$H_k=\1\{\sigma(k)\text{'s prefix is in }C_{\mathsf{mem}_P}\text{ at }\alpha_k\}$.
The link and the decode batch drain $X_k$ and $D_k$ at the rates \eqref{L1:eq:rates}.
\end{lemma}

\begin{proof}
A request joins the admission queue of $\mathsf{mem}_P$ at $a_k$ and can be admitted
only at its head, after $\alpha_{k-1}$, and only once the guard of \textsc{Admit}
holds, which is \eqref{L1:eq:admitP}. After its admission it executes \textsc{Run} at
once, so it enters $Q_{\mathsf{prefill}}$ in admission order, which is arrival order;
by \eqref{L1:eq:rates} it is worked at rate $1$ exactly when it is first, that is from
the \textsc{Finish} of its predecessor or its own entry, whichever is later, until
$S_k$ is worked off, which is \eqref{L1:eq:lindley}. After \textsc{Finish} at
$\mathsf{decode}$ the session frees its memory, branches, runs $Z_k$ at the external
stage (rate $1$, no waiting), executes \textsc{Turn} and \textsc{Loop}, and queues
again, all in zero time except the tool, which is \eqref{L1:eq:feedback}. The
statement about $H_k$ is the first assignment of \textsc{Admit}.
\end{proof}

\Cref{L1:eq:feedback} is the closed loop: request arrival times depend on completion
times, which depend on the congestion. The request rate is nevertheless fixed by the
session rate. If the system is stable, every session leaves after its $J_i$ turns
(\cref{L1:ex:program}), and
\begin{equation}\label{L1:eq:turnrate}
  \lambda=\Lambda\,\E[J]=\frac{\Lambda}{1-p}.
\end{equation}

\paragraph{The stages as queues.} Each stage kind is a queue studied in one later
lecture. A serial stage is a single first-in first-out server, and
\eqref{L1:eq:lindley} is its recursion: the M/G/1 queue of \cref{sec:L3} when arrivals
are Poisson and the closed system of \cref{sec:L4} when a finite number of sessions
circulate. A shared stage is a processor-sharing station (\cref{sec:L2}), and an
external stage is a delay station (\cref{sec:L4}). A pool is a queue too. If every
context needs about $c$ bytes, the guard of \textsc{Admit} reads $|R_P(t)|<m_P$ with
$m_P=\lfloor M_P/c\rfloor$ \emph{slots}: a multi-server queue $\cdot$/G/$m_P$
(\cref{sec:L2}) whose holding time is the time from \textsc{Admit} to \textsc{Free},
\[
  e_k-\alpha_k={\underbrace{(s_k-\alpha_k)}_{\text{wait for prefill}}}+S_k
  +{\underbrace{(e_k-f_k)}_{\text{transfer}}},
\]
and likewise $m_D=\lfloor M_D/c\rfloor$ slots held for $d_k-\alpha'_k$ at the decode
instance. With unequal contexts a request needs several slots, and the queue has no
closed form.

\paragraph{The cache chooses the prefill work.} In classical queueing theory the
service-time law is given. Here, by \eqref{L1:eq:work} and \textsc{Admit}, the cache
chooses it: the eviction order $E$ decides $H_k$, and
\begin{equation}\label{L1:eq:ES}
  \E[S]=\E\bigl[P(n,K)\bigr]+\E\bigl[(1-H)\,\Delta S(K)\bigr].
\end{equation}
Only the second term depends on the policy, and it depends on it through which
sessions miss, weighted by $\Delta S(K)$, which grows as $K^2$. The central question of
\cref{sec:L3} is how a performance measure changes when the distribution of $S$
changes; \eqref{L1:eq:ES} says the eviction policy is what changes it.

\paragraph{Performance.} The first token is produced when the prefill ends, so
\begin{equation}\label{L1:eq:ttft}
  \mathrm{TTFT}_k=f_k-a_k
  ={\underbrace{(\alpha_k-a_k)}_{W^M_{P,k}}}
  +{\underbrace{(s_k-\alpha_k)}_{W_{P,k}}}+S_k,
\end{equation}
the wait for KV memory at the prefill instance, the wait for the prefill instance, and
the prefill work. The transfer and the decode come after the first token and do not
enter it; they set when the rest of the answer arrives, and they reach the TTFT only
through memory, because a request keeps its prefill allocation until $e_k$. The
objective of the scheduler is the time average of the number of requests in the
system,
\[
  L=\lim_{t\to\infty}\frac1t\int_0^t N(u)\dd u,\qquad N(u)=\#\{k:a_k\le u<d_k\},
\]
which by Little's law (\cref{thm:little}) equals $\lambda W$ with $W$ the average of
$d_k-a_k$ over requests and $\lambda$ from \eqref{L1:eq:turnrate}. No stage can be
stable unless its load is below its capacity:
\begin{equation}\label{L1:eq:stable}
  \rho_P=\lambda\,\E[S]<1,\qquad \rho_L=\lambda\,\E[X]<1,\qquad
  \lambda\,\E[D]<\max_{m\le m_D}\varphi(m).
\end{equation}
By \eqref{L1:eq:ES}, the first condition already depends on the eviction policy.

\paragraph{What the model leaves out.} Work is a continuous quantity rather than whole
tokens and memory blocks; an instance's scheduler steps are replaced by the continuous
flow \eqref{L1:eq:rates}; the allocation $c_k$ ignores the $o_k$ tokens a decode adds,
and a running decode that needs a block when none is free is preempted and later
recomputed, which the step stage with $\kw{growing}$ expresses but the replica of
\cref{L1:ex:program} does not use; and by \eqref{L1:eq:feedback}
request arrivals are not Poisson. \Cref{sec:L5} returns to these gaps once the tools
are in place.

\section{The exponential distribution and memorylessness}

\begin{definition}[Exponential distribution]\label{L1:def:exp}
A random variable $X\ge0$ with $\Prob(X>t)=e^{-\mu t}$ for $t\ge0$ has the
\term{exponential distribution} with \term{rate} $\mu>0$, written
$X\sim\mathrm{Exp}(\mu)$. Then $\E[X]=1/\mu$ and $\E[X^2]=2/\mu^2$.
\end{definition}

\begin{proposition}[Memorylessness \cite{ross1996}]\label{prop:memoryless}
If $X\sim\mathrm{Exp}(\mu)$, then for all $s,t\ge0$
\begin{equation}\label{L1:eq:memoryless}
  \Prob(X>s+t\mid X>s)=\Prob(X>t).
\end{equation}
Conversely, if a random variable with $0<X<\infty$ a.s.\ satisfies
\eqref{L1:eq:memoryless}, then $X\sim\mathrm{Exp}(\mu)$ for some $\mu\in(0,\infty)$.
\end{proposition}

\begin{proof}
The forward direction is $e^{-\mu(s+t)}/e^{-\mu s}=e^{-\mu t}$.
Converse: let $G(t)=\Prob(X>t)$. Then \eqref{L1:eq:memoryless} reads
$G(s+t)=G(s)G(t)$. $G$ is nonincreasing and right-continuous, and
$G(0)=\Prob(X>0)=1$. By induction $G(nt)=G(t)^n$, hence $G(1)=G(1/n)^n$ and
$G(m/n)=G(1)^{m/n}$.
If $G(1)=0$, then $G(1/n)=0$ for every $n$, which contradicts
$G(1/n)\to G(0)=1$ by right-continuity. If $G(1)=1$, then $G(r)=1$ for every
rational $r$, and by monotonicity $G\equiv1$, that is $X=\infty$ a.s., contradicting
the assumption. Hence $G(1)=e^{-\mu}$ for some $\mu\in(0,\infty)$, and
$G(r)=e^{-\mu r}$ on the rationals. For arbitrary $t$, take rationals $r\downarrow t$;
right-continuity gives $G(t)=e^{-\mu t}$.
\end{proof}

Memorylessness says that the remaining service time of a customer who has already
received $s$ units of service has the same distribution as the service time of a new
customer. This is why, in a system with exponential service, the number in system
$N(t)$ alone determines the future, and $N(t)$ is a continuous-time Markov chain
(\cref{sec:L2}). The following fact is also used often.

\begin{lemma}[Minimum of exponentials \cite{ross1996}]\label{L1:lem:min}
If $X_i\sim\mathrm{Exp}(\mu_i)$, $i=1,\dots,n$, are independent, then
$\min_iX_i\sim\mathrm{Exp}(\sum_i\mu_i)$.
\end{lemma}

\begin{proof}
$\Prob(\min_iX_i>t)=\prod_i\Prob(X_i>t)=e^{-(\sum_i\mu_i)t}$.
\end{proof}

\section{The Poisson process}

\begin{definition}[Counting process and Poisson process]\label{def:poisson}
A \term{counting process} $A=(A(t))_{t\ge0}$ is a nondecreasing, right-continuous,
$\N$-valued process with $A(0)=0$ and jumps of size one only.
A counting process $A$ is a \term{Poisson process} with \term{rate} $\lambda>0$ if
\begin{enumerate}
\item (\term{independent increments}) for disjoint intervals
  $(s_1,t_1],\dots,(s_n,t_n]$ the increments $A(t_j)-A(s_j)$ are independent;
\item (\term{stationary Poisson increments}) for all $s\ge0$, $t>0$,
  $A(s+t)-A(s)\sim\mathrm{Poisson}(\lambda t)$, that is
  $\Prob(A(s+t)-A(s)=n)=e^{-\lambda t}(\lambda t)^n/n!$.
\end{enumerate}
$S_n=\inf\{t:A(t)\ge n\}$ is the $n$-th arrival time, and $T_n=S_n-S_{n-1}$
($S_0=0$) is the $n$-th \term{interarrival time}.
\end{definition}

\begin{theorem}[Exponential interarrival times \cite[\S2.2]{ross1996}]\label{L1:thm:interarrival}
$A$ is a Poisson process with rate $\lambda$ if and only if the interarrival times
$T_1,T_2,\dots$ are independent and each is $\mathrm{Exp}(\lambda)$.
\end{theorem}

\begin{proof}[Proof ($\Rightarrow$)]
For $0=t_0<t_1<\dots<t_n$ and small $h>0$ consider the event
$E_h=\{S_j\in(t_j,t_j+h],\ j=1,\dots,n\}$. Let $F_h$ be the event that there is no
arrival in $(t_{j-1}+h,t_j]$ and exactly one in each $(t_j,t_j+h]$. Then
$F_h\subset E_h$, and $E_h\setminus F_h$ is contained in the event that some interval
of length $h$ has two or more arrivals, which has probability $O(h^{n+1})$. By
independent and Poisson increments,
\[
  \Prob(F_h)=e^{-\lambda(t_1+(t_2-t_1-h)+\cdots+(t_n-t_{n-1}-h))}\,(\lambda h e^{-\lambda h})^n
  =\lambda^n h^n e^{-\lambda t_n}\,(1+O(h)).
\]
Hence $(S_1,\dots,S_n)$ has density $\lambda^ne^{-\lambda s_n}$ on
$\{0<s_1<\dots<s_n\}$. The linear map $T_j=S_j-S_{j-1}$ has Jacobian one and
$s_n=\sum_jT_j$, so $(T_1,\dots,T_n)$ has density $\prod_{j=1}^n\lambda e^{-\lambda t_j}$,
that is, independent $\mathrm{Exp}(\lambda)$. The converse is a standard argument
using that $S_n$ is Gamma distributed together with memorylessness, and we omit it
\cite[\S2.2]{ross1996}.
\end{proof}

\begin{proposition}[Superposition and thinning \cite{kingman1993}]\label{prop:superpose}
\begin{enumerate}[label=(\roman*)]
\item (\term{Superposition}) If $A_1,A_2$ are independent Poisson processes with rates
  $\lambda_1,\lambda_2$, then $A_1+A_2$ is a Poisson process with rate
  $\lambda_1+\lambda_2$.
\item (\term{Thinning}) If each point of a Poisson process with rate $\lambda$ is kept
  with probability $p$ and discarded with probability $1-p$, independently of
  everything else, then the kept points and the discarded points are Poisson
  processes with rates $p\lambda$ and $(1-p)\lambda$, and they are independent.
\end{enumerate}
\end{proposition}

\begin{proof}
(i) Independent increments follow from the independent increments of the two
processes and their mutual independence. For the increment over one interval, the
generating function of independent Poisson variables is
$\E[z^{X_1+X_2}]=e^{\lambda_1t(z-1)}e^{\lambda_2t(z-1)}=e^{(\lambda_1+\lambda_2)t(z-1)}$,
so the sum is $\mathrm{Poisson}((\lambda_1+\lambda_2)t)$.

(ii) Let $X$ be the original increment over an interval of length $t$, $X_1$ the number
kept, and $X_2=X-X_1$ the number discarded. Given $X=j+k$, the conditional
probability that $X_1=j$ is binomial, $\binom{j+k}{j}p^j(1-p)^k$, so
\[
  \Prob(X_1=j,X_2=k)=e^{-\lambda t}\frac{(\lambda t)^{j+k}}{(j+k)!}\binom{j+k}{j}p^j(1-p)^k
  =e^{-p\lambda t}\frac{(p\lambda t)^j}{j!}\cdot e^{-(1-p)\lambda t}\frac{((1-p)\lambda t)^k}{k!}.
\]
Thus on one interval $X_1,X_2$ are independent Poisson variables. On disjoint
intervals the original increments and the coin flips are all independent, so all
increments of the two processes are independent.
\end{proof}

\serving{If sessions arrive as a Poisson process of rate $\Lambda$ and each turn is
routed independently to one of $J$ prefill instances, then by
\cref{prop:superpose}(ii) each instance also sees Poisson arrivals. The turns within a
session, however, go around a closed loop (prefill $\to$ link $\to$ decode $\to$
tool), so the turn arrivals are not exactly Poisson. The closed models of
\cref{sec:L4} are the tool for this gap.}

\section{PASTA: Poisson arrivals see time averages}

The distribution of the state seen by arriving customers generally differs from the
time-average distribution (\cref{L1:ex:lookahead}(b) below gives an example). With
Poisson arrivals this discrepancy disappears.

\begin{definition}[Lack of anticipation]\label{L1:def:lookahead}
Given a state process $X=(X(t))_{t\ge0}$ and an arrival process $A$, we say that $A$
has \term{lack of anticipation} with respect to $X$ if for every $t$ the future
increments $\{A(t+s)-A(t):s\ge0\}$ are independent of the past
$\{X(u):u\le t\}\cup\{A(u):u\le t\}$.
\end{definition}

In a queueing system this assumption says that future arrivals are independent of the
system state so far. Note the direction: the state may (and does) depend on past
arrivals; what is excluded is that future arrivals react to the state. Poisson arrivals
from outside satisfy it naturally. Arrivals in a closed system (\cref{sec:L4}) do not.

\begin{example}[One process with lack of anticipation, one without]\label{L1:ex:lookahead}
\emph{(a) Lack of anticipation: the M/M/1 queue.} Customers come from a large outside
population as a Poisson process, independently of the service times. The queue length
$X(t)$ is a function of the arrivals before $t$ and the service times, and by
independent increments the arrivals after $t$ are independent of both. An arriving
customer is like a raindrop falling on the queue: the rain does not know how long the
queue is. So the arrival instants are ``random inspection times'', and by
\cref{thm:pasta} the fraction of arrivals that find the server busy equals the fraction
of time it is busy, $\rho$.

\emph{(b) No lack of anticipation: a single agent session.} One session alternates
between a tool call (mean $Z$) and a turn at a server (mean $1/\mu$). Its
next turn can arrive only after its previous turn has left the server. If $X(t)=1$
(the turn is being served), then no arrival can happen until the service ends; if
$X(t)=0$, an arrival may come at any moment. Future arrivals depend on the current
state, so the condition fails. The consequence is visible: every arrival finds the
server idle, while the server is busy a fraction $(1/\mu)/(Z+1/\mu)>0$ of the time.
Arrivals do not see time averages. With $N$ sessions the same effect is weaker but
present, and \cref{sec:L4} replaces PASTA by the arrival theorem for exactly this
reason.
\end{example}

\begin{remark}[Why study PASTA if sessions violate it?]\label{L1:rem:why-pasta}
Example~\ref{L1:ex:lookahead}(b) shows that a single agent session does not have lack
of anticipation. PASTA is still the right starting point, for three reasons.
\begin{enumerate}
\item \emph{The prefill model rests on it.} The prefill instance is modelled as an
  M/G/1 queue, and the proof of the Pollaczek--Khinchine formula (\cref{thm:pk}) uses
  PASTA. Session feedback makes this an approximation, to be tested against the
  system. One cannot judge an approximation without knowing what it approximates.
\item \emph{Many sessions look Poisson.} Each session reacts to the state, but the
  superposition of many independent sessions, each contributing a small share of the
  arrivals, is close to a Poisson process; this is the Palm--Khinchine theorem
  \cite{khinchine1960}. New sessions themselves arrive as a Poisson process in the
  model, so the M/G/$\infty$ and Campbell calculations of \cref{sec:L4} use Poisson
  arrivals without approximation.
\item \emph{The exact closed-system answer is a corrected PASTA.} The arrival theorem
  (\cref{thm:arrival}) says that an arriving turn sees the time averages of the system
  with one session fewer, and it reduces to PASTA as the number of sessions grows
  (\cref{L4:rem:arrival-limit}). For exponential work, \cref{prop:finite}(ii) shows
  that the open formula built on PASTA overstates the wait at equal utilisation, so
  PASTA gives a conservative estimate whose error the finite-population model
  quantifies.
\end{enumerate}
\end{remark}

\begin{theorem}[PASTA \cite{wolff1982}]\label{thm:pasta}
Let $A$ be a Poisson process with rate $\lambda$, $X$ a state process with left
limits, and $B$ a set of states, and suppose $A$ has lack of anticipation with respect
to $X$. If $U(t)=\1\{X(t^-)\in B\}$ is left-continuous, then for every $T>0$
\begin{equation}\label{L1:eq:pasta}
  \E\Bigl[\sum_{k:\,S_k\le T}\1\{X(S_k^-)\in B\}\Bigr]=\lambda\,\E\Bigl[\int_0^T\1\{X(t)\in B\}\dd t\Bigr].
\end{equation}
Consequently, if the limits of the two averages
$\frac1T\int_0^T\1\{X(t)\in B\}\dd t$ and
$\frac{1}{A(T)}\sum_{k\le A(T)}\1\{X(S_k^-)\in B\}$ exist with probability one, they
are equal: the fraction of arrivals that find the system in $B$ equals the fraction of
time it spends in $B$.
\end{theorem}

\begin{proof}
The left side is $\int_0^TU(t)\dd A(t)$. For the partition $t_j=jT/n$ of $[0,T]$,
define the step function $U_n(t)=\sum_{j}U(t_j)\1\{t\in(t_j,t_{j+1}]\}$. $U(t_j)$ is
determined by the past up to time $t_j$, and by lack of anticipation the increment
$A(t_{j+1})-A(t_j)$ is independent of that past, so
\[
  \E\Bigl[\int_0^TU_n\dd A\Bigr]=\sum_j\E[U(t_j)]\,\E[A(t_{j+1})-A(t_j)]
  =\lambda\sum_j\E[U(t_j)]\,(t_{j+1}-t_j)=\lambda\,\E\Bigl[\int_0^TU_n(t)\dd t\Bigr].
\]
Since $U$ is left-continuous, $U_n(t)\to U(t)$ for every $t$. The integrands on the
two sides are bounded by $A(T)$ and $T$ respectively, and $\E[A(T)]=\lambda T<\infty$,
so by dominated convergence $\E[\int_0^TU\dd A]=\lambda\E[\int_0^TU\dd t]$ as
$n\to\infty$. Finally, the set of $t$ with $X(t^-)\ne X(t)$ is countable, so replacing
$U(t)$ by $\1\{X(t)\in B\}$ on the right does not change the integral.
The statement about limits follows from a sample-path version of
\eqref{L1:eq:pasta} (a law of large numbers applied to a martingale), for which we
cite \cite{wolff1982}.
\end{proof}

In \cref{sec:L3}, PASTA becomes a key step of the Pollaczek--Khinchine formula in the
form ``the mean work an arriving turn finds waiting equals the time-average work
waiting''.

\section{Little's law}

\begin{theorem}[Little's law \cite{little1961,stidham1974}]\label{thm:little}
For a sequence of customers $(a_k,d_k)$ with $a_k\le d_k<\infty$ and $a_k$
nondecreasing, suppose the two limits
\[
  \lambda=\lim_{t\to\infty}\frac{A(t)}{t}\in(0,\infty),\qquad
  W=\lim_{n\to\infty}\frac1n\sum_{k=1}^nW_k<\infty
\]
exist. Then $L=\lim_{t\to\infty}\frac1t\int_0^tN(s)\dd s$ exists and
\begin{equation}\label{L1:eq:little}
  L=\lambda W .
\end{equation}
\end{theorem}

\begin{proof}
Since $N(s)=\sum_k\1\{a_k\le s<d_k\}$, we have
$\int_0^tN(s)\dd s=\sum_k\bigl|[a_k,d_k)\cap[0,t]\bigr|$.
Each term is at most $W_k$, is zero if $a_k>t$, and equals $W_k$ exactly if
$d_k\le t$. Therefore
\begin{equation}\label{L1:eq:little-sandwich}
  \sum_{k:\,d_k\le t}W_k\;\le\;\int_0^tN(s)\dd s\;\le\;\sum_{k:\,a_k\le t}W_k=\sum_{k=1}^{A(t)}W_k .
\end{equation}
Upper bound: since $A(t)\to\infty$,
$\frac1t\sum_{k\le A(t)}W_k=\frac{A(t)}{t}\cdot\frac{1}{A(t)}\sum_{k\le A(t)}W_k\to\lambda W$.

Lower bound: with $\bar W_n=\frac1n\sum_{k\le n}W_k$ we have
$\frac{W_n}{n}=\bar W_n-\frac{n-1}{n}\bar W_{n-1}\to W-W=0$. Also $A(a_n)\ge n$ and
$A(a_n^-)\le n-1$, so $a_n/n\to1/\lambda$, and therefore $W_n/a_n\to0$.
For any $\varepsilon>0$ there is $k_0$ such that $d_k=a_k+W_k\le(1+\varepsilon)a_k$ for
$k\ge k_0$. Then for $t\ge\max_{k<k_0}d_k$, every customer with
$a_k\le t/(1+\varepsilon)$ has $d_k\le t$, so
\[
  \frac1t\int_0^tN(s)\dd s\;\ge\;\frac1t\sum_{k=1}^{A(t/(1+\varepsilon))}W_k
  \longrightarrow\frac{\lambda W}{1+\varepsilon}.
\]
Letting $\varepsilon\downarrow0$, the liminf is at least $\lambda W$. Together with
the upper bound this gives \eqref{L1:eq:little}.
\end{proof}

The proof uses no distributional assumption and no service discipline. It is a
deterministic theorem that holds on every sample path on which the limits exist. So
depending on how we draw the boundary of ``the system'', we obtain several identities.

\begin{corollary}[Utilisation law \cite{kleinrock1975}]\label{cor:busy}
Under the assumptions of \cref{thm:little}, consider a single-server system and
suppose the average of the service times $\frac1n\sum_{k\le n}S_k$ converges to
$\E[S]$.
\begin{enumerate}[label=(\roman*)]
\item The fraction of time the server is busy is $\rho=\lambda\E[S]$. In particular
  $\rho\le1$ in a stable system.
\item For the waiting room, $L_q=\lambda W_q$, and under FIFO $L=L_q+\rho$.
\end{enumerate}
\end{corollary}

\begin{proof}
(i) Take ``the server'' alone as the system. Customer $k$ enters it when its service
starts and leaves $S_k$ later. Every customer is eventually served, so the arrival
rate of this subsystem is also $\lambda$, and the customer average of its sojourn times
is $\E[S]$. The number in this subsystem is 1 when the server is busy and 0 otherwise,
so its time average is the fraction of time the server is busy. By \cref{thm:little}
this equals $\lambda\E[S]=\rho$, and being a fraction it is at most one.
(ii) Taking the waiting room alone as the system, the sojourn times are $W_{q,k}$, so
$L_q=\lambda W_q$. Taking time averages in $N=N_q+(\text{number in service})$ gives
$L=L_q+\rho$.
\end{proof}

\begin{example}\label{L1:ex:little}
Suppose $\lambda=2$ turns per second arrive at a prefill instance and the mean prefill
work is $\E[S]=0.3$ seconds. By \cref{cor:busy} the prefill server is busy a fraction
$\rho_P=0.6$ of the time. If the observed mean number of turns waiting for it is
$L_q=0.9$, then $W_q=L_q/\lambda=0.45$ seconds, the mean TTFT (with no wait for memory)
is $W_q+\E[S]=0.75$ seconds, and the mean number of turns at the prefill server is
$L=\lambda\cdot0.75=1.5$.
\end{example}

The scheduler's objective is the mean number of turns in the stations, for example
$L=L_P+L_D$ for the prefill server and the decode batch. By \cref{thm:little} this
equals the total delay accumulated per unit time, $\lambda\cdot(\text{mean sojourn
time})$, so ``reduce the mean number of turns'' and ``reduce the mean delay'' are the
same goal. The ``price of a miss'' of \cref{sec:L3} is defined precisely as a change in
$L_P$.

\section{Exercises}

\begin{exercise}\label{L1:exr:race}
For independent $X_1\sim\mathrm{Exp}(\mu_1)$ and $X_2\sim\mathrm{Exp}(\mu_2)$, show that
$\Prob(X_1<X_2)=\mu_1/(\mu_1+\mu_2)$, and that the event $\{X_1<X_2\}$ is independent
of $\min(X_1,X_2)$.
\emph{Hint:} integrate $\Prob(X_1<X_2,\ \min>t)$ against the density of $X_1$.
\end{exercise}

\begin{exercise}\label{L1:exr:delay}
Each time a session makes a tool call, it stays away from the server for a mean of $Z$
seconds. If tool calls start at rate $\lambda$, show that the time-average number of
sessions in a tool call is $\lambda Z$. Does this hold whatever the distribution of the
tool-call duration?
\emph{Hint:} take the tool-call stage alone as ``the system'' and apply
\cref{thm:little}.
\end{exercise}

\begin{exercise}\label{L1:exr:nocache}
Remove $\kw{cache}\ \kappa T$ from the program of \cref{L1:ex:program}. Show from
\cref{L1:def:semantics} that every turn then has $H_k=0$, and compute $\E[S]$ from
\eqref{L1:eq:ES}. Which of the conditions \eqref{L1:eq:stable} becomes harder to meet?
\end{exercise}

\begin{exercise}\label{L1:exr:colocated}
In a \emph{colocated} replica, prefill and decode run on one instance with one pool of
$M$ bytes, and no transfer is needed. Write its \textsc{Route} program with one step
stage. The instance of the paper gives each iteration's budget first to the decoding
turns and the rest to the oldest prefill; the step stage serves its jobs in their
order of entry. Show that the two orders coincide when turns enter in the order they
are admitted and a prefill may take all the budget left: after every iteration every
job in its decode phase precedes every job in its prefill phase. Give a two-job
example showing that a cap on a single prefill's tokens per iteration breaks this.
\emph{Hint:} a prefill that does not finish in an iteration leaves no budget behind it.
\end{exercise}
